\section{Amazon Virtual Private Cloud}

Amazon VPC là miền mạng logic cô lập trong AWS. CIDR xác định dải địa chỉ;
subnet chia dải đó theo Availability Zone; route table quyết định target cho
từng prefix; Security Group (SG) lọc traffic stateful tại network interface
\cite{aws-vpc-guide}. Trong thí nghiệm, bốn CIDR không chồng lấn là điều kiện
để route liên VPC có đích xác định.

\section{VPC Peering}

VPC Peering nối hai VPC để traffic private đi trực tiếp giữa chúng. Kết nối
không có tính bắc cầu: A peering với B và B peering với C không làm A tự động
truy cập C qua B \cite{aws-peering-guide}. Vì vậy, full mesh của $n$ VPC cần:

\begin{equation}
N_{\mathrm{peering}}=\frac{n(n-1)}{2}.
\end{equation}

Với $n=4$, hệ thống cần sáu connection: AB, AC, AD, BC, BD và CD. Mỗi route
table VPC vẫn cần ba route có hướng tới ba CIDR còn lại, tổng cộng 12 route
liên VPC trong cấu hình này.

\section{AWS Transit Gateway}

TGW cung cấp mô hình hub-and-spoke. Mỗi VPC tạo một attachment tới hub; TGW
route table điều khiển việc chuyển tiếp giữa các attachment
\cite{aws-tgw-overview,aws-tgw-routing}. Với bốn VPC, đồ án dùng một TGW và
bốn VPC attachment. Mô hình này giảm số quan hệ kết nối logic khi số VPC tăng,
nhưng không làm mất yêu cầu cấu hình route đi và route về.

\section{So sánh kiến trúc Peering và TGW}

\begin{table}[H]
\centering
\caption{Đặc tính kiến trúc, chưa phải kết quả hiệu năng}
\label{tab:architecture-comparison}
\begin{tabularx}{\textwidth}{p{0.22\textwidth}XX}
\toprule
\textbf{Tiêu chí} & \textbf{VPC Peering} & \textbf{Transit Gateway} \\
\midrule
Topology & Point-to-point; full mesh khi mọi VPC cần liên lạc & Hub-and-spoke \\
Tính bắc cầu & Không hỗ trợ & Chuyển tiếp qua hub theo TGW route table \\
Quản trị & Phân tán theo từng cặp và route table VPC & Tập trung attachment và TGW route table \\
Mở rộng kết nối & Số connection tăng theo $n(n-1)/2$ & Số attachment tăng tuyến tính theo $n$; quota áp dụng \cite{aws-tgw-quotas} \\
Chi phí cơ bản & Không tính phí theo giờ cho connection; data transfer áp dụng theo tài liệu AWS \cite{aws-peering-pricing} & Có phí attachment theo giờ và data processing \cite{aws-tgw-pricing} \\
\bottomrule
\end{tabularx}
\end{table}

Bảng \ref{tab:architecture-comparison} chỉ mô tả thiết kế. RTT và goodput phải
được xác định từ thực nghiệm, không suy ra trực tiếp từ số connection.

\section{EC2 networking và ENA}

EC2 network performance phụ thuộc loại instance, kích thước, loại traffic và
chính sách allowance. Tài liệu AWS phân biệt baseline bandwidth với khả năng
burst ở một số instance; mức công bố không đồng nghĩa một TCP flow luôn đạt
mức đó \cite{aws-ec2-bandwidth}. ENA xuất các counter như
\nolinkurl{bw_in_allowance_exceeded} và
\nolinkurl{pps_allowance_exceeded}. Counter tăng cho biết traffic đã gặp giới
hạn tương ứng tại interface; không được cộng các counter để suy ra packet loss
\cite{aws-ena-metrics}.

Trong fan-in, ba sender hội tụ vào B. Ngay cả khi đường kết nối còn dư năng
lực, giới hạn bandwidth, PPS, CPU hoặc xử lý TCP tại B vẫn có thể chi phối tổng
goodput. Vì vậy CPU và ENA là bằng chứng hỗ trợ xác định bottleneck endpoint,
không phải metric phụ.

\section{TCP và các metric}

RTT là thời gian khứ hồi của một request và response. Throughput là tốc độ dữ
liệu truyền qua hệ thống; receiver goodput trong báo cáo là số bit payload mà
phía nhận ghi nhận chia cho thời lượng nhận. Với số byte $B$ và thời gian
$T$, giá trị dùng trong báo cáo là:

\begin{equation}
G=\frac{8B}{T\times10^9}\quad\text{(Gbps)}.
\end{equation}

TCP retransmission là số segment TCP truyền lại do iperf3 báo cáo. Đại lượng
này không đồng nhất với phần trăm packet loss. Một single flow còn chịu ảnh
hưởng của congestion window, RTT, socket processing và giới hạn endpoint;
nhiều concurrent flow có thể phân chia tài nguyên khác với một flow.

\section{Mô hình fan-in}

Kịch bản fan-in có ba source cùng gửi tới một receiver:

\begin{verbatim}
A -----5201---\
C -----5202----> B
D -----5203---/
\end{verbatim}

Ba listener tách cổng để tránh tranh chấp một tiến trình server và giữ từng
JSON độc lập. Tổng fan-in hiện tại là tổng ba goodput tính trên cửa sổ riêng
của từng source. Do wrapper khởi chạy lệch dưới một giây, đây là tổng xấp xỉ,
không phải phép tích phân byte trên một cửa sổ giao nhau tuyệt đối.

\section{Metric thống kê}

Với $m=5$ lượt và giá trị $x_i$, mean và sample standard deviation là:

\begin{align}
\bar{x}&=\frac{1}{m}\sum_{i=1}^{m}x_i,\\
s&=\sqrt{\frac{\sum_{i=1}^{m}(x_i-\bar{x})^2}{m-1}}.
\end{align}

Coefficient of variation là $CV=100s/\bar{x}$. Chênh lệch tương đối dùng
Peering làm mốc được ghi rõ tại từng bảng. Mức chia đều ba flow trong một lượt
dùng Jain fairness index \cite{jain-fairness}:

\begin{equation}
J=\frac{(\sum_{i=1}^{3}g_i)^2}{3\sum_{i=1}^{3}g_i^2}.
\end{equation}

$J$ gần 1 chỉ mô tả mức chia đều trong các cửa sổ đã ghi; không chứng minh
max-min fairness hoặc năng lực tối đa của dịch vụ.

\section{Phân vị và chỉ số đúng hạn}

Mean có thể che giấu một nhóm request rất chậm. Báo cáo vì vậy dùng phân vị
P95 theo quy tắc nearest rank. Với $n$ độ trễ đã sắp tăng dần,

\begin{equation}
k=\lceil0{,}95n\rceil,\qquad P95=x_{(k)}.
\end{equation}

P95 trong nghiên cứu ứng dụng chỉ tính trên request hoàn thành thành công.
Để không loại request lỗi hoặc trễ khỏi mẫu số chất lượng dịch vụ, nhóm định
nghĩa SLI đúng hạn:

\begin{equation}
\widehat p=\frac{N_{\mathrm{good}}}{N_{\mathrm{scheduled}}},
\qquad SLI_{\%}=100\widehat p.
\end{equation}

$N_{\mathrm{good}}$ là request thành công và kết thúc không quá 50 ms kể từ
lịch gửi; $N_{\mathrm{scheduled}}$ gồm toàn bộ request theo lịch. Ngưỡng 50 ms
là lựa chọn của thí nghiệm, không phải SLA của AWS. P50/P95/P99 và SLI/SLO được
liên hệ với bài giảng NT531 về Internet Performance và QoS
\cite{nt531-qos-lecture}.

\section{Công suất phục vụ danh nghĩa và giới hạn đồng thời}

Mỗi worker có $M$ slots. Một lời gọi chiếm một slot trong thời gian cấu hình
$s=0{,}010$ s. Với tải theo lịch $\lambda_{\mathrm{cfg}}$, công suất và mức tải
danh nghĩa được dùng để chọn điều kiện thử:

\begin{equation}
\mu_{\mathrm{nom}}=\frac{M}{s},\qquad
\rho_{\mathrm{nom}}=\frac{\lambda_{\mathrm{cfg}}s}{M}.
\end{equation}

Với 8 slots, $\mu_{\mathrm{nom}}=800$ lời gọi/s; 800 request/s làm
$\rho_{\mathrm{nom}}=1$. Với 16 slots, công suất danh nghĩa là 1.600 lời gọi/s
và $\rho_{\mathrm{nom}}=0{,}5$. Đây là mô hình chọn tải, không phải giới hạn AWS,
không phải phần trăm CPU và không giả định quá trình đến Poisson.

\section{Định luật Little cho hàng đợi quan sát}

Bài giảng Arrival Processes and Queuing Systems trình bày quan hệ
$L=\lambda W$ và ví dụ thread pool \cite{nt531-queue-lecture}. Nhóm áp dụng cho
phần request đang chờ:

\begin{equation}
\overline L_{q,C}=\frac{1}{T_C}\sum_{i=1}^{N}W_{q,i}
=\underbrace{\frac{N}{T_C}}_{\lambda_C}
 \underbrace{\frac{1}{N}\sum_{i=1}^{N}W_{q,i}}_{\overline W_q}.
\end{equation}

$T_C$ bao phủ từ lần vào đầu tiên đến lần thoát cuối cùng của cùng cohort, kể
cả đuôi xử lý. Vì vậy $\lambda_C=N/T_C$ có thể khác rate cấu hình. Hai vế được
tính từ cùng log để kiểm tra nhất quán, không phải hai mẫu thống kê độc lập.

\section{Quan sát phân tán}

OpenTelemetry mô hình hóa một request bằng trace và các span; Collector nhận
telemetry, Tempo lưu trace, còn Grafana dùng để truy vấn và hiển thị
\cite{otel-docs}. Pipeline này giúp nhìn vị trí queue/service trong đường gọi
A $\rightarrow$ B $\rightarrow$ C,D. Trace smoke được dùng làm minh họa cơ chế;
benchmark chính tắt full tracing để tránh thay đổi workload cần so sánh.
